% #####################################################################
\chapter{Levantamento e análise de requisitos}
\label{textuais/cap4}


\section{Descrição do problema e escopo}
% DESCRIÇÃO: resuma o problema (irrigação sem controle desperdiça água;
% soluções comerciais são caras e voltadas a outras escalas) e delimite o
% escopo: um único dispositivo (esp32-01), irrigação por gotejamento em
% pequenas áreas, acompanhamento por painel web. Diga também o que ficou
% fora do escopo (ex.: vários dispositivos, fertirrigação).

\section{Perfis de usuário}
% DESCRIÇÃO: descreva quem usa o sistema e o que cada perfil pode fazer:
%  - Administrador: configura parâmetros, culturas e latitude, aciona a
%    irrigação, muda o modo, autoriza o modo automático e cadastra
%    usuários;
%  - Visualizador: acompanha status, sensores e histórico, sem alterar
%    nada;
%  - Dispositivo (ator do sistema): envia leituras e recebe comandos.

\section{Requisitos funcionais}
% DESCRIÇÃO: tabela com identificador (RF01, RF02...), descrição e perfil.
% Exemplos tirados dos repositórios:
%  - monitorar umidade do solo, temperatura e umidade do ar;
%  - alternar entre os modos manual e automático;
%  - autorizar ou bloquear a irrigação automática;
%  - acionar a irrigação manual com volume em litros;
%  - cadastrar culturas e seus Kc por fase;
%  - configurar os parâmetros agronômicos e a latitude;
%  - exibir histórico em gráficos (24 h, 7 e 30 dias) e eventos de
%    irrigação;
%  - autenticar usuários e permitir que o administrador cadastre novos.

\section{Requisitos não funcionais}
% DESCRIÇÃO: tabela com identificador (RNF01...), descrição e categoria
% (segurança, desempenho, usabilidade, tecnologia). Exemplos:
%  - autenticação de usuários pelo Firebase Authentication;
%  - autenticação do dispositivo por chave (X-Device-Key);
%  - armazenamento em nuvem (Firestore);
%  - atualização do painel a cada 5 s (status) e 10 s (telemetria);
%  - interface web responsiva;
%  - comunicação do dispositivo por Wi-Fi e HTTPS.

\section{Regras de negócio}
% DESCRIÇÃO: as regras que o sistema segue para decidir. Ex.:
%  - a irrigação automática só ocorre se o modo for automático, se
%    estiver autorizada e se a umidade do solo estiver abaixo do limiar;
%  - o Kc é definido pela cultura e pela fase escolhidas, não digitado;
%  - o volume aplicado depende dos parâmetros agronômicos;
%  - somente administradores alteram configurações e acionam comandos.

\section{Modelagem do sistema}
% DESCRIÇÃO: diagramas que ajudam a banca a entender o sistema antes da
% implementação. Sugestões:
%  - diagrama de casos de uso (perfis x funcionalidades);
%  - diagrama de sequência de um fluxo principal (ex.: comando de
%    irrigação manual saindo do painel até o ESP32).

\begin{figure}[htbp]
  \centering
  % \includegraphics[width=\textwidth]{figuras/casos-de-uso}
  \caption{Diagrama de casos de uso do sistema}
  \label{fig:casos-de-uso}
  \fonte{Elaborado pelo autor.}
\end{figure}

